%% Introduction au problème %%

Le principal frein à l'accessibilité de l'ensemble des collections de l'INA est l'absence de compréhension par ces publics, de l'organisation des dites collections. En ignorant comment a été réalisé la collecte, la conservation, la description et l'indexation de ces documents, il est difficile de savoir comment effectuer sa recherche. Dans l'ensemble, les utilisateurs des bases de données de l'INA, avant la médiation des documentalistes, n'ont pas vraiment conscience de ce que conserve l'INA. Le périmètre de leurs fonds sont flou pour un public externe. Ils vont par exemple chercher à l'INA des images de la première guerre mondiale quand ces images sont conservées à l'ECPAD (Établissement de communication et de production audiovisuelle de la Défense). En conséquence, les utilisateurs ont encore moins conscience des politiques documentaire. Il est difficile de savoir dans quel champs chercher. Pour une émission diffusée avant 1975, il vaut mieux ne pas utiliser de mots clés mais d'un autre côté, en effectuant une recherche plein-texte, les résultats sont très bruyants. De plus, quand on ne connais pas la base de données comme les connaissent les documentalistes, il est aisé de faire un contre-sens, de faire une recherche, d'avoir un certain nombre de résultats et s'y tenir, sans savoir qu'en utilisant un autre champs ou un autre mots-clé, les résultats auraient été bien plus nombreux\footnote{(Voir entretien d'une documentaliste de l'INA sur l'expertise du documentaliste~\ref{ann:transcriptions},\nameref{ann:transcriptions}), p.~\pageref{ann:transcriptions}.}. 

%% Application à chaque public %%

C'est alors que les questions se dirigent vers les documentalistes, détenteurs de ce savoir, soit pour se décharger de la recherche documentaire comme le font les professionnels de l’audiovisuel dont la recherche n'est pas le domaine et qui ne souhaitent pas en faire souffrir leur créations, soit pour acquérir les connaissances nécessaires pour effectuer leurs recherches. L'expertise des documentalistes est d'une grande valeurs. C'est pourquoi une de leurs missions consiste à le concrétiser en outils facilitant la compréhension de ces collections comme le sont la création de corpus et l'écriture de billet de blog pour le blog \enquote{Hypothèse}. Malgré tout, les documentalistes ne sont pas omniscients, leur connaissance a des limites. Elle ne peut pas s'étendre aux nombreux cas particuliers qu'ils n'ont jamais étudiés ou à la connaissance d'indexations utilisée avant qu'ils ne prennent leur poste. 

Ce premier frein à la recherche dans les bases de données des collections de l'INAthèque est également implicite dans les autres problèmes que nous allons soulever. Chaque frein à la pleine accessibilité des collections existent en partie parce que les utilisateurs n'en ont pas forcément conscience. 